\documentclass[11pt]{article}

\usepackage[margin=1in]{geometry}
\usepackage{amsmath,amssymb}
\usepackage{array}
\usepackage{booktabs}
\usepackage{enumitem}
\usepackage[hidelinks]{hyperref}

\setlist[itemize]{leftmargin=1.4em}
\setlist[enumerate]{leftmargin=1.6em}

\title{SYSC 5108 Exam Study\\Assignment 4, Question 6}
\author{Jesse Levine}
\date{}

\begin{document}
\maketitle

\section*{Question}

Given the two-dimensional dataset
\[
\begin{array}{c|ccccc}
    x_1 & 2 & 0 & -3 & 3 & -2\\
    x_2 & -2 & 4 & -1 & -3 & 2
\end{array}
\]
find:
\begin{enumerate}[label=(\alph*)]
    \item the one-dimensional PCA,
    \item the encoded points,
    \item the reconstructed points,
    \item the percentage of variance explained.
\end{enumerate}

\section*{Course and Textbook Cross-References}

\begin{itemize}
    \item \textbf{Chapter 7 slides, PCA:} PCA finds low-dimensional components
    that preserve as much variance as possible.
    \item \textbf{Chapter 7 slides, Linear PCA and Interpretation:} the principal
    directions come from the covariance structure of the data.
    \item \textbf{Goodfellow, Bengio, and Courville, \emph{Deep Learning}, PCA derivation:}
    the encoding is a linear projection onto principal directions, and the
    reconstruction is obtained from the projected code.
\end{itemize}

\section*{Step 1: Write the Data Points}

The five samples are
\[
    x^{(1)}=(2,-2),\quad
    x^{(2)}=(0,4),\quad
    x^{(3)}=(-3,-1),\quad
    x^{(4)}=(3,-3),\quad
    x^{(5)}=(-2,2).
\]

\section*{Step 2: Compute the Sample Mean}

The mean is
\[
    \mu
    =
    \frac{1}{5}
    \sum_{i=1}^5 x^{(i)}
    =
    \frac{1}{5}
    \begin{bmatrix}
        2+0-3+3-2\\
        -2+4-1-3+2
    \end{bmatrix}
    =
    \begin{bmatrix}
        0\\
        0
    \end{bmatrix}.
\]

So the data are already centered.

\section*{Step 3: Form the Data Matrix}

Using columns as data points,
\[
    X
    =
    \begin{bmatrix}
        2 & 0 & -3 & 3 & -2\\
        -2 & 4 & -1 & -3 & 2
    \end{bmatrix}.
\]

\section*{Step 4: Compute the Covariance Matrix}

Because the mean is zero, the sample covariance matrix used in the submitted
solution is
\[
    \Sigma = \frac{1}{5}XX^\top.
\]

First compute \(XX^\top\):
\[
    XX^\top
    =
    \begin{bmatrix}
        2 & 0 & -3 & 3 & -2\\
        -2 & 4 & -1 & -3 & 2
    \end{bmatrix}
    \begin{bmatrix}
        2 & -2\\
        0 & 4\\
        -3 & -1\\
        3 & -3\\
        -2 & 2
    \end{bmatrix}
    =
    \begin{bmatrix}
        26 & -14\\
        -14 & 34
    \end{bmatrix}.
\]

Therefore
\[
    \Sigma
    =
    \frac{1}{5}
    \begin{bmatrix}
        26 & -14\\
        -14 & 34
    \end{bmatrix}
    =
    \begin{bmatrix}
        5.2 & -2.8\\
        -2.8 & 6.8
    \end{bmatrix}.
\]

\section*{Step 5: Find the Principal Eigenvalue and Eigenvector}

The principal component is the eigenvector of \(\Sigma\) corresponding to the
largest eigenvalue.

Solve
\[
    \det(\Sigma-\lambda I)=0:
\]
\[
    \det
    \begin{bmatrix}
        5.2-\lambda & -2.8\\
        -2.8 & 6.8-\lambda
    \end{bmatrix}
    =0.
\]

So
\[
    (5.2-\lambda)(6.8-\lambda)-(-2.8)^2=0,
\]
\[
    \lambda^2-12\lambda+27.52=0.
\]

The two eigenvalues are
\[
    \lambda_1 \approx 8.9120,
    \qquad
    \lambda_2 \approx 3.0880.
\]

The largest eigenvalue is \(\lambda_1\), so the one-dimensional PCA direction is
its unit eigenvector:
\[
    u_1
    =
    \begin{bmatrix}
        0.6022\\
        -0.7983
    \end{bmatrix}.
\]

\section*{Answer to Part (a): One-Dimensional PCA}

\[
    \boxed{
    u_1=
    \begin{bmatrix}
        0.6022\\
        -0.7983
    \end{bmatrix}
    }
\]

Note: the sign can be flipped and still represent the same principal direction.

\section*{Step 6: Encode the Data Points}

For one-dimensional PCA, each encoded coordinate is the projection
\[
    z_i = u_1^\top x^{(i)}.
\]

Compute each one:
\[
    z_1 = [0.6022\;\; -0.7983]
    \begin{bmatrix}
        2\\
        -2
    \end{bmatrix}
    = 2.8011,
\]
\[
    z_2 = [0.6022\;\; -0.7983]
    \begin{bmatrix}
        0\\
        4
    \end{bmatrix}
    = -3.1934,
\]
\[
    z_3 = [0.6022\;\; -0.7983]
    \begin{bmatrix}
        -3\\
        -1
    \end{bmatrix}
    = -1.0082,
\]
\[
    z_4 = [0.6022\;\; -0.7983]
    \begin{bmatrix}
        3\\
        -3
    \end{bmatrix}
    = 4.2016,
\]
\[
    z_5 = [0.6022\;\; -0.7983]
    \begin{bmatrix}
        -2\\
        2
    \end{bmatrix}
    = -2.8011.
\]

\section*{Answer to Part (b): Encoded Points}

\[
    \boxed{
    z=
    \begin{bmatrix}
        2.8011 &
        -3.1934 &
        -1.0082 &
        4.2016 &
        -2.8011
    \end{bmatrix}
    }
\]

\section*{Step 7: Reconstruct the Data}

The PCA reconstruction formula is
\[
    \hat{x}^{(i)} = \mu + u_1 z_i.
\]
Since \(\mu=0\),
\[
    \hat{x}^{(i)}=u_1 z_i.
\]

Thus
\[
    \hat{x}^{(1)}
    =
    2.8011
    \begin{bmatrix}
        0.6022\\
        -0.7983
    \end{bmatrix}
    \approx
    \begin{bmatrix}
        1.6868\\
        -2.2362
    \end{bmatrix},
\]
\[
    \hat{x}^{(2)}
    =
    -3.1934
    \begin{bmatrix}
        0.6022\\
        -0.7983
    \end{bmatrix}
    \approx
    \begin{bmatrix}
        -1.9230\\
        2.5494
    \end{bmatrix},
\]
\[
    \hat{x}^{(3)}
    =
    -1.0082
    \begin{bmatrix}
        0.6022\\
        -0.7983
    \end{bmatrix}
    \approx
    \begin{bmatrix}
        -0.6072\\
        0.8049
    \end{bmatrix},
\]
\[
    \hat{x}^{(4)}
    =
    4.2016
    \begin{bmatrix}
        0.6022\\
        -0.7983
    \end{bmatrix}
    \approx
    \begin{bmatrix}
        2.5302\\
        -3.3544
    \end{bmatrix},
\]
\[
    \hat{x}^{(5)}
    =
    -2.8011
    \begin{bmatrix}
        0.6022\\
        -0.7983
    \end{bmatrix}
    \approx
    \begin{bmatrix}
        -1.6868\\
        2.2362
    \end{bmatrix}.
\]

\section*{Answer to Part (c): Reconstructed Points}

\[
    \boxed{
    \hat{x}^{(1)}=(1.6868,-2.2362),\quad
    \hat{x}^{(2)}=(-1.9230,2.5494),\quad
    \hat{x}^{(3)}=(-0.6072,0.8049),}
\]
\[
    \boxed{
    \hat{x}^{(4)}=(2.5302,-3.3544),\quad
    \hat{x}^{(5)}=(-1.6868,2.2362).}
\]

\section*{Step 8: Percentage of Variance Explained}

For PCA, the variance explained by the first component is
\[
    \frac{\lambda_1}{\lambda_1+\lambda_2}.
\]

Substitute the eigenvalues:
\[
    \frac{8.9120}{8.9120+3.0880}
    =
    \frac{8.9120}{12}
    \approx 0.7427.
\]

So the explained variance percentage is
\[
    0.7427 \times 100\% \approx 74.27\%.
\]

\section*{Answer to Part (d): Variance Explained}

\[
    \boxed{\text{Explained variance ratio} \approx 0.7427}
\]
\[
    \boxed{\text{Explained variance percentage} \approx 74.27\%}
\]

\section*{Final Answer Summary}

\begin{itemize}
    \item One-dimensional PCA direction:
    \[
        u_1=
        \begin{bmatrix}
            0.6022\\
            -0.7983
        \end{bmatrix}
    \]
    \item Encoded points:
    \[
        z=[2.8011,\ -3.1934,\ -1.0082,\ 4.2016,\ -2.8011]
    \]
    \item Reconstructed points:
    \[
        (1.6868,-2.2362),\ (-1.9230,2.5494),\ (-0.6072,0.8049),\ (2.5302,-3.3544),\ (-1.6868,2.2362)
    \]
    \item Variance explained:
    \[
        74.27\%
    \]
\end{itemize}

\section*{Exam Shortcut}

For a small PCA question:
\begin{enumerate}[label=\arabic*.]
    \item Center the data.
    \item Compute the covariance matrix.
    \item Take the eigenvector with the largest eigenvalue.
    \item Encode with projection \(z_i=u_1^\top x_i\).
    \item Reconstruct with \(\hat{x}_i=\mu+u_1 z_i\).
    \item Compute explained variance as \(\lambda_1/\sum_j \lambda_j\).
\end{enumerate}

\section*{Sources Used}

\begin{itemize}
    \item Assignment 4 PDF, Question 6.
    \item Local submitted Assignment 4 PDF checked for consistency of the numerical results.
    \item Chapter 7 slides on PCA, linear PCA, and the covariance interpretation.
    \item Goodfellow, Bengio, and Courville, \emph{Deep Learning}, PCA derivation via projection and reconstruction.
\end{itemize}

\end{document}
